\subsection{A signed-only full display}
\Needspace{9\baselineskip}
\begin{proposition}[Full occupancy does not make the criteria agree]
\label{prop:full-four-distinction}
There is a full-support $16\times4$ display of length $43=U(16,4)$
on a maximum singleton base which satisfies $\SG$ but not $\RW$.
\end{proposition}
\begin{proof}
Number cells by $a=4r+c$, with indices starting at zero. The singleton
row supports are (writing $01$ for $\{0,1\}$, and similarly)
\[
 01,02,03,12,13,23,2,2,3,3,3,3,3,3,3,3.
\]
Take the following unordered pairs:
\begin{align*}
 &(2,3),(5,60),(7,49),(9,62),(10,57),(12,21),(15,29),\\
 &(16,50),(18,61),(24,31),(25,48),(27,28),(32,38),(33,40),\\
 &(34,36),(37,46),(41,56),(44,45),(20,54),(42,52),(53,58).
\end{align*}
The 22 singletons form a maximum base by Lemma~\ref{lem:cap}.
Every other cell occurs exactly once, so the length is $43=U(16,4)$.

Here the distinction is entirely in orthogonality: every selected pair
normalizes recursively. Packet F16 gives a 31-step normalization
sequence and a checker. After these identifications, form the
graph on the $\binom{43}{2}=903$ products of distinct selected classes.
Line products are grounded; rectangle equations reverse signs along edges.
A direct inventory gives exactly three components without a zero ground,
of sizes $5,13,3$. Each contains a triangle. For example, writing
$A=[53,58]$, $B=[42,52]$, $C=[41,56]$ for three paired classes
and $[a]$ for a singleton class, one triangle has vertices
\[
 [0]A,\qquad [1]B,\qquad [2,3]C.
\]
Its rectangle witnesses are $(0,13;0,1)$, $(0,10;1,2)$ and
$(0,14;0,2)$, where $(r,s;c,d)$ denotes rows $r,s$ and columns $c,d$. Each product denotes the inner product of the corresponding
class vectors. The signed rule grounds all three odd components. The
recursive rules cannot enter them. The F16 reader reconstructs the entire
product graph and checks all three triangles, proving the two outcomes.
\end{proof}

This example does not separate the optimized values, even on this base.
Replacing $(20,54),(42,52),(53,58)$ by
$(20,42),(52,58),(53,54)$ recovers the recursive length-$43$
construction of \cite[Figure~2]{CC}. Thus full occupancy does not justify keeping the pairing unchanged
when converting a signed display to a recursive one.
